\section{A society of agents that learn from each other}
\label{sec:model}

\subsection{Agents, issues and opinions}

A society of $N$ agents discusses an agenda of $P$ issues. An issue is a vector
$\hat{x}\in\mathbb{R}^{K}$: the embedding of an assertion in the $K$ dimensions
that are relevant to the question at hand. All agents embed issues identically
and differ only in what they make of them, so the model separates the
representation of an issue from the opinion an agent holds about it.

Each agent carries a single-layer perceptron with weights $w_I\in\mathbb{R}^{K}$
and states the opinion $\sigma_I = \mathrm{sign}(w_I\cdot\hat{x})\in\{\pm 1\}$.
This is the least structure that lets two agents agree on some issues and
disagree on others; a model in which each agent is a single spin cannot do that,
and the possibility of partial agreement turns out to drive the collective
behaviour we report. Agents also belong to one of two classes, $A$ or $B$,
recorded as $\kappa_I=\pm1$: a label that is visible to others and carries no
information about any issue.

An agent's state has two sectors. The \emph{ideological} sector is a Gaussian
belief over its own weights, with mean $\hat{w}_I$ and covariance $C_I$. The
\emph{affective} sector is a Gaussian belief, for every other agent $J$, over how
unreliable $J$ is as a source: mean $\mu_{J|I}$ (the \emph{distrust} $I$ assigns
to $J$) and variance $V_{J|I}$. Writing the joint belief as a product of
Gaussians,
%
\begin{equation}
  Q(w_I, z_I \mid \lambda_I) \;=\; \mathcal{N}(\hat{w}_I, C_I)
  \prod_{J\neq I}\mathcal{N}(\mu_{J|I}, V_{J|I}),
  \label{eq:belief}
\end{equation}
%
the dynamical variables of the society are
$\lambda_I=\{\hat{w}_I, C_I, \mu_{J|I}, V_{J|I}\}$. The means fix what an agent
thinks; the (co)variances fix how firmly, and thereby how fast it can be moved.

\subsection{One interaction}

At each step an issue $\hat{x}$, an emitter $e$ and a receiver $r$ are drawn
uniformly. The emitter states $\sigma_e$; the receiver treats it as a noisy
report of an unknown true label, where the noise level is precisely what its
distrust of the emitter encodes. A Bayesian step on \eqref{eq:belief} leaves the
exponential family, and projecting the posterior back onto a Gaussian by maximum
entropy --- relative to the prior, subject to the first two moments --- keeps the
belief in the form \eqref{eq:belief} and yields a closed update. This is the
entropic-dynamics construction, and it is what makes the algorithm optimal in the
small-group (student--teacher) setting rather than merely plausible.
Appendix~\ref{app:derivation} carries the derivation, including the two modelling
choices that keep the algebra Gaussian: taking the probability that the emitter
conveys the wrong label to be $\Phi(z)$ with $z$ Gaussian, and letting an
inexperienced agent be correspondingly unsure it has parsed the issue as the
emitter did.

Two scaled fields carry everything the update depends on:
%
\begin{equation}
  \hw \;=\; \frac{\hat{w}_r\cdot\hat{x}\,\sigma_e}{\gamma_C},
  \qquad
  \hmu \;=\; \frac{\mu_{e|r}}{\gamma_V},
  \qquad
  \gamma_C=\sqrt{1+\hat{x}\cdot C_r\hat{x}},\quad \gamma_V=\sqrt{1+V_{e|r}} .
  \label{eq:fields}
\end{equation}
%
We call $\hw$ the \emph{opinion field}: it is positive when the receiver already
agrees with the message and large when it holds that opinion confidently. We
call $\hmu$ the \emph{distrust field}: positive when the receiver distrusts the
emitter. In these variables the evidence for the message is
%
\begin{equation}
  Z \;=\; \Phi(\hw) + \Phi(\hmu) - 2\,\Phi(\hw)\Phi(\hmu),
  \label{eq:evidence}
\end{equation}
%
which is small in two quadrants --- agreeing with a distrusted emitter, and
disagreeing with a trusted one --- and large in the other two. Low evidence is
surprise, and surprise is what drives learning. The update is a gradient step in
$\log Z$ in both sectors,
%
\begin{align}
  \hat{w}_r &\leftarrow \hat{w}_r + \frac{\Fw}{\gamma_C}\,\sigma_e\,C_r\hat{x},
  &
  C_r &\leftarrow C_r + \frac{\FC}{\gamma_C^{2}}\,(C_r\hat{x})(C_r\hat{x})^{\!\top},
  \label{eq:update-ideo}\\
  \mu_{e|r} &\leftarrow \mu_{e|r} + \frac{\Fmu}{\gamma_V}\,V_{e|r},
  &
  V_{e|r} &\leftarrow V_{e|r} + \frac{\FV}{\gamma_V^{2}}\,V_{e|r}^{2},
  \label{eq:update-aff}
\end{align}
%
with the four modulation functions given by the derivatives of $\log Z$:
%
\begin{align}
  \Fw &= \bigl(1-2\Phi(\hmu)\bigr)\frac{g(\hw)}{Z},
  &\Fmu &= \bigl(1-2\Phi(\hw)\bigr)\frac{g(\hmu)}{Z},
  \nonumber\\
  \FC &= -\Fw(\Fw+\hw),
  &\FV &= -\Fmu(\Fmu+\hmu),
  \label{eq:modulation}
\end{align}
%
where $g$ is the standard normal density. The first update in
\eqref{eq:update-ideo} is a Hebbian step $\sigma_e\hat{x}$ with a tensorial,
self-annealing learning rate; $\FC$ and $\FV$ shrink the two uncertainties, so
an agent's opinions and its judgements of others both harden with experience.

Three properties of \eqref{eq:modulation} do the work in what follows, and none of
them was put in by hand. \textbf{Trust gates learning and can invert it}: the
prefactor $1-2\Phi(\hmu)$ in $\Fw$ is negative when the emitter is distrusted, so
a distrusted emitter is not ignored but \emph{anti-learned} from, generating
effectively antiferromagnetic couplings. (A comparable reversal is documented in
humans, where a norm advertised by a distrusted source can push behaviour the
other way, \citealp{keizer2011reversal}.) \textbf{Agreement builds trust}:
symmetrically, $1-2\Phi(\hw)$ in $\Fmu$ is negative when the receiver agrees, so
whether $\mu$ rises or falls is decided by the \emph{sign of the perceived
agreement} --- the hinge that \S\ref{sec:discrimination} turns.
\textbf{Only one sector yields at a time}: the sectors are mirror images,
$\Fw(x,y)=\Fmu(y,x)$ and $\FC(x,y)=\FV(y,x)$, and the surprise is absorbed by
whichever field is smaller in magnitude, with a crossover at $\hmu=\hw$
(Figure~\ref{fig:modulation}). Meeting a distrusted agent who agrees with you
costs you either your distrust or your opinion, depending on which you hold less
firmly --- blame attribution emerging from the inference problem rather than being
assumed.

\begin{figure}[t]
  \centering
  \begin{subfigure}{0.49\linewidth}
    \includegraphics[width=\linewidth]{figures/modulation_contours}
    \caption{The four modulation functions.}
  \end{subfigure}\hfill
  \begin{subfigure}{0.43\linewidth}
    \includegraphics[width=\linewidth]{figures/modulation_slices}
    \caption{Which sector absorbs the surprise.}
  \end{subfigure}
  \caption{\textbf{Learning is driven by dissonance.}
  (a) $\Fmu, \Fw$ (top) and $\FV, \FC$ (bottom) over the $(\hw,\hmu)$ plane; the
  green line is $\hmu=\hw$, and reflection across it exchanges the sectors. Changes
  concentrate in the dissonant quadrants: agree-and-distrust (upper right),
  disagree-and-trust (lower left). (b) Cuts at fixed $\hw$: below the crossover at
  $\hmu=\hw+D$ the affective sector moves (solid), above it the ideological sector
  does (dashed) and the receiver unlearns. The discrimination field shifts that
  crossover, which is the whole mechanism studied here.}
  \label{fig:modulation}
\end{figure}
